% GENERATED FILE. Edit canonical lessons, then run npm run generate:books.
% canonical-source-hash: fnv1a64:a306317d31d628b7
% canonical-lessons: ES-C584-ventaja, ES-C584-desventaja, ES-C584-inconveniente, ES-C584-positivo, ES-C584-negativo

\chapter{Pros y contras}
\label{ch:es-pros-y-contras}
\hlchaptermodality{584}

\begin{chapteropening}
\textbf{By the end of this chapter:} \emph{I can weigh the advantages and the drawbacks of a choice.}

Complete the last lesson of chapter 584: I can weigh the advantages and the drawbacks of a choice.
\end{chapteropening}

\section[la ventaja]{la ventaja — the advantage}

\label{lesson:ES-C584-ventaja}



\begin{quote}
Before the new one: say the Spanish for the factor, then the Spanish for to generate.
\end{quote}

\subsection*{You'll want to know: la ventaja}
\textbf{la ventaja} — \textquotedblleft{}the advantage\textquotedblright{}.

\begin{quote}
Vivir cerca tiene muchas ventajas.
\end{quote}

\begin{quote}
\emph{Living nearby has many advantages.}
\end{quote}

\emph{Sacar ventaja} — to gain an advantage, to get ahead.

\begin{cousinweb}
\emph{ventaja} was borrowed from French \emph{avantage} 'a being in front', built on \emph{avant}, from Latin \emph{ab ante} 'in front' (EtymDB).
\end{cousinweb}

\begin{grammarlens}[title={What you've built}]
One more word for telling, explaining and describing at B1.
\end{grammarlens}

\subsection*{Your turn}
\begin{itemize}
\raggedright
  \item \emph{Say it:} \emph{la ventaja}
  \item \emph{Say it:} \emph{la ventaja}, once more
  \item \emph{Say it:} a sentence of your own with \emph{la ventaja}
  \item \emph{Recall:} say the Spanish for the factor, then the Spanish for to generate, then say \emph{la ventaja} again
\end{itemize}

\begin{tcolorbox}[breakable,colback=teal!4,colframe=teal!35!black,title={Before you move on}]
What does la ventaja mean? (\textquotedblleft{}The advantage\textquotedblright{}.) Say it once more.
\end{tcolorbox}
\section[la desventaja]{la desventaja — the disadvantage}

\label{lesson:ES-C584-desventaja}



\begin{quote}
Before the new one: say the Spanish for to generate, then the Spanish for the advantage.
\end{quote}

\subsection*{You'll want to know: la desventaja}
\textbf{la desventaja} — \textquotedblleft{}the disadvantage\textquotedblright{}.

\begin{quote}
La desventaja es el precio.
\end{quote}

\begin{quote}
\emph{The disadvantage is the price.}
\end{quote}

The opposite of \emph{la ventaja}.

\begin{cousinweb}
\emph{desventaja} is built on \emph{ventaja} with the prefix \emph{des-}; \emph{ventaja} was borrowed from French \emph{avantage}, from Latin \emph{ab ante} 'in front' (EtymDB).
\end{cousinweb}

\begin{grammarlens}[title={What you've built}]
One more word for telling, explaining and describing at B1.
\end{grammarlens}

\subsection*{Your turn}
\begin{itemize}
\raggedright
  \item \emph{Say it:} \emph{la desventaja}
  \item \emph{Say it:} \emph{la desventaja}, once more
  \item \emph{Say it:} a sentence of your own with \emph{la desventaja}
  \item \emph{Recall:} say the Spanish for to generate, then the Spanish for the advantage, then say \emph{la desventaja} again
\end{itemize}

\begin{tcolorbox}[breakable,colback=teal!4,colframe=teal!35!black,title={Before you move on}]
What does la desventaja mean? (\textquotedblleft{}The disadvantage\textquotedblright{}.) Say it once more.
\end{tcolorbox}
\section[el inconveniente]{el inconveniente — the drawback, the problem; inconvenient}

\label{lesson:ES-C584-inconveniente}



\begin{quote}
Before the new one: say the Spanish for the advantage, then the Spanish for the disadvantage.
\end{quote}

\subsection*{You'll want to know: el inconveniente}
\textbf{el inconveniente} — \textquotedblleft{}the drawback, the problem; inconvenient\textquotedblright{}.

\begin{quote}
No hay ningún inconveniente.
\end{quote}

\begin{quote}
\emph{There is no problem at all.}
\end{quote}

\emph{No tengo inconveniente} — I do not mind.

\begin{cousinweb}
\emph{inconveniente} is a learned borrowing from Latin \emph{inconveniens} 'unsuitable', built on \emph{convenire} 'to come together, to suit' with \emph{in-} 'not' (EtymDB).
\end{cousinweb}

\begin{grammarlens}[title={What you've built}]
One more word for telling, explaining and describing at B1.
\end{grammarlens}

\subsection*{Your turn}
\begin{itemize}
\raggedright
  \item \emph{Say it:} \emph{el inconveniente}
  \item \emph{Say it:} \emph{el inconveniente}, once more
  \item \emph{Say it:} a sentence of your own with \emph{el inconveniente}
  \item \emph{Recall:} say the Spanish for the advantage, then the Spanish for the disadvantage, then say \emph{el inconveniente} again
\end{itemize}

\begin{tcolorbox}[breakable,colback=teal!4,colframe=teal!35!black,title={Before you move on}]
What does el inconveniente mean? (\textquotedblleft{}The drawback, the problem; inconvenient\textquotedblright{}.) Say it once more.
\end{tcolorbox}
\section[positivo]{positivo — positive}

\label{lesson:ES-C584-positivo}



\begin{quote}
Before the new one: say the Spanish for the disadvantage, then the Spanish for the drawback.
\end{quote}

\subsection*{You'll want to know: positivo}
\textbf{positivo} — \textquotedblleft{}positive\textquotedblright{}.

\begin{quote}
Tiene una actitud positiva.
\end{quote}

\begin{quote}
\emph{She has a positive attitude.}
\end{quote}

\emph{Lo positivo} — the good side.

\begin{cousinweb}
\emph{positivo} is a learned borrowing from Latin \emph{positivus} 'set down, settled', built on \emph{ponere} 'to put' (EtymDB).
\end{cousinweb}

\begin{grammarlens}[title={What you've built}]
One more word for telling, explaining and describing at B1.
\end{grammarlens}

\subsection*{Your turn}
\begin{itemize}
\raggedright
  \item \emph{Say it:} \emph{positivo}
  \item \emph{Say it:} \emph{positivo}, once more
  \item \emph{Say it:} a sentence of your own with \emph{positivo}
  \item \emph{Recall:} say the Spanish for the disadvantage, then the Spanish for the drawback, then say \emph{positivo} again
\end{itemize}

\begin{tcolorbox}[breakable,colback=teal!4,colframe=teal!35!black,title={Before you move on}]
What does positivo mean? (\textquotedblleft{}Positive\textquotedblright{}.) Say it once more.
\end{tcolorbox}
\section[negativo]{negativo — negative}

\label{lesson:ES-C584-negativo}



\begin{quote}
Before the new one: say the Spanish for the drawback, then the Spanish for positive.
\end{quote}

\subsection*{You'll want to know: negativo}
\textbf{negativo} — \textquotedblleft{}negative\textquotedblright{}.

\begin{quote}
No seas tan negativo.
\end{quote}

\begin{quote}
\emph{Do not be so negative.}
\end{quote}

\emph{Lo negativo} — the bad side.

\begin{cousinweb}
\emph{negativo} is a learned borrowing from Latin \emph{negativus}, built on \emph{negare} 'to deny' (EtymDB).
\end{cousinweb}

\begin{grammarlens}[title={What you've built}]
One more word for telling, explaining and describing at B1.
\end{grammarlens}

\subsection*{Your turn}
\begin{itemize}
\raggedright
  \item \emph{Say it:} \emph{negativo}
  \item \emph{Say it:} \emph{negativo}, once more
  \item \emph{Say it:} a sentence of your own with \emph{negativo}
  \item \emph{Recall:} say the Spanish for the drawback, then the Spanish for positive, then say \emph{negativo} again
\end{itemize}

\begin{tcolorbox}[breakable,colback=teal!4,colframe=teal!35!black,title={Before you move on}]
What does negativo mean? (\textquotedblleft{}Negative\textquotedblright{}.) Say it once more.
\end{tcolorbox}
